% ===================================================================
%  TABEL Nr. 10 — De zee. — De haven.
%  Traduction 2026 d'apres texto/io, controlee sur texto/fr et
%  texto/es. Consigne : voir texto/nl/10-tabel-01.tex.
%
%  Le francais decale sa numerotation d'un cran a partir de l'alinea
%  8 ; c'est une coquille de l'atelier francais, et la colonne suit
%  l'ido, qui compte juste.
% ===================================================================

%%K t10-apar-1 apar
\VUsaut{4.0mm}
\VUtitre{15.0pt}{60}{TABEL Nr. 10}
\VUsaut{3.0mm}
\VUpk{11.4pt}{40}{De zee. --- De haven.}
\VUsaut{3.0mm}

%%K t10-01-1 p
1. --- In de zomer, terwijl sommigen rust en koelte op een berg zoeken, gaan anderen liever
naar de kust. Zo deden wij vorig jaar, want wij houden veel van de \VUgras{Oceaan}
\textsuperscript{(1)}.

%%K t10-02-1 p
2. --- Wat mij betreft, ik voel een groot genoegen bij het aanschouwen van die onmetelijke
watervlakte die zich tot aan de \VUgras{horizon} \textsuperscript{(2)} uitstrekt, en bij het
zien vertrekken voor een lange
\VUgras{overtocht} van de enorme \VUgras{stoomschepen}
\textsuperscript{(3)} waarop de reizigers inschepen die naar Amerika gaan.

%%K t10-03-1 p
3. --- Daarom kiest mijn vader, die mijn neiging kent, om de vakantie door te brengen altijd
een heel rustig strand, dat bij een van onze grote
\VUgras{oorlogshavens} ligt.

%%K t10-03-2 p
Tijdens ons laatste verblijf kwam het \VUgras{eskader} achter de enorme
\VUgras{dijk} \textsuperscript{(4)} voor anker gaan die de \VUgras{oorlogshaven}
\textsuperscript{(5)} afsluit en beschermt, en het lukte mij naar wens de verschillende
\VUgras{oorlogsschepen} te bewonderen, de kleine
\VUgras{onderzeeboot} \textsuperscript{(10)} die duikt en onder het water
verdwijnt en ook het enorme \VUgras{pantserschip} \textsuperscript{(9)} dat tot nu toe bijna
alleen de aanvallen van de \VUgras{torpedoboot} \textsuperscript{(8)} vreesde.

%%K t10-tit-2 sub
\VUsaut{3.0mm}
\VUpk{10.2pt}{40}{De kleine schepen.}
\VUsaut{3.0mm}

%%K t10-04-1 p
4. --- Deze laatste wist al heel \VUgras{behendig} het zwakke punt van het dikke metalen
pantser te vinden om daar de gevaarlijke \VUgras{torpedo} te steken, waarvan de ontploffing de
ondergang van het schip veroorzaakt, maar toch was zij minder te duchten dan de onderzeeboot,
want de torpedobootjagers zetten haar gemakkelijk na.

%%K t10-05-1 p
5. --- Ik kon ook de snelle gepantserde \VUgras{kruisers} \textsuperscript{(6)} zien, bijna
even onkwetsbaar als het echte pantserschip, verscheidene
\VUgras{transportschepen} \textsuperscript{(12)} drie of vier
\VUgras{kanonneerboten} \textsuperscript{(7)} en ten slotte een
\VUgras{fregat} \textsuperscript{(11)}. \VUgras{Het schouwspel van die vloot,
wanneer zij manoeuvreert om het anker te lichten is het boeiendst.}

%%K t10-06-1 p
6. --- Wat een verschil in bouw tussen die grote staalmassa's en de sierlijke
\VUgras{driemasters} of de bevallige \VUgras{zeilschepen}
\textsuperscript{(13)} die nu nog in de koopvaardijvloot gebruikt worden en die ik met alle
zeilen de haven zag binnenvaren, toen ik op de \VUgras{kade} \textsuperscript{(14)} wandelde.
Weliswaar verloren deze veel van hun voordelen
\VUgras{wanneer} de \VUgras{wind} tegen was. Zij moesten alle zeilen strijken om
het bassin weer binnen te varen en een \VUgras{sleepboot} \textsuperscript{(16)} was nodig om
ze te halen en ze, als eenvoudige
\VUgras{vrachtschuiten} \textsuperscript{(17)}, naar de kade te brengen.

%%K t10-tit-3 sub
\VUsaut{3.0mm}
\VUpk{10.2pt}{40}{De handelshaven.}
\VUsaut{3.0mm}

%%K t10-07-1 p
7. --- Wat boeiend is in de haven waarover ik spreek, is dat zij niet alleen een oorlogshaven
bevat, kunstmatig door reusachtige werken aangelegd, maar ook een wonderbare
\VUgras{handelshaven} die in de \VUgras{monding} van een grote stroom ligt.

%%K t10-08-1 p
8. --- De landtong die de twee bassins scheidt is versterkt en verdedigd door een reeks
\VUgras{batterijen} en \VUgras{bastions} die in zee vooruitspringen. Men heeft een bijzondere
vergunning nodig om op de \VUgras{wallen} \textsuperscript{(15)} te wandelen. Ik kreeg die
gemakkelijk, door mijn oom, die ingenieur van de \VUgras{scheepswerven} \textsuperscript{(18)}
is en ik ging de schepen bezoeken die men onder ruime loodsen bouwt.

%%K t10-09-1 p
9. --- Ik woonde zelfs de tewaterlating van een pantserschip bij en ik herinner mij het
aangrijpende ogenblik dat de hele menigte voelde toen de reusachtige massa, aan zichzelf
overgelaten, over haar wiegvormige stellage begon te glijden en op het water van de stroom
kwam drijven.

%%K t10-10-1 p
10. --- Om naar de scheepswerven te gaan, steekt men het
\VUgras{exercitieterrein} \textsuperscript{(19)} over waar de soldaten van het
garnizoen elke dag komen oefenen, men gaat over een ijzeren \VUgras{brugje}
\textsuperscript{(20)}, dat het paradeplein met het \VUgras{arsenaal} \textsuperscript{(21)}
verbindt.

%%K t10-tit-4 sub
\VUsaut{3.0mm}
\VUpk{10.2pt}{40}{Het arsenaal en de dokken.}
\VUsaut{3.0mm}

%%K t10-11-1 p
11. --- Met mijn oom bezocht ik dat grote gebouw vol \VUgras{kanonnen} \textsuperscript{(22)},
\VUgras{kogels} \textsuperscript{(23)} en
\VUgras{granaten}. Ook zag ik de \VUgras{metaalfabriek}
\textsuperscript{(24)} waarin de \VUgras{ketels} \textsuperscript{(25)} van de stoomschepen
gemaakt worden.

%%K t10-12-1 p
12. --- Vlakbij zijn de \VUgras{dokken} \textsuperscript{(26)} --- reparatiedokken --- en de
heel merkwaardige \VUgras{drijvende dokken} \textsuperscript{(27)} waarin ik van heel
dichtbij, op een te herstellen schip, de \VUgras{schroef} \textsuperscript{(28)}, de
\VUgras{kiel} \textsuperscript{(29)} en het \VUgras{roer} \textsuperscript{(30)} van een boot
kon zien.

%%K t10-13-1 p
13. --- De handelshaven is even bestuderenswaard als de oorlogshaven, en ik bracht veel uren
door met gapen op de kaden waar de \VUgras{raderboten} \textsuperscript{(31)} gemeerd liggen
die de stroom weer opvaren, en met te kijken naar het werken van de \VUgras{masthijskranen}
\textsuperscript{(32)}, die, als veren, enorme lasten optillen.

%%K t10-tit-5 sub
\VUsaut{4.0mm}
\VUpk{10.2pt}{40}{De loods.}
\VUsaut{3.0mm}

%%K t10-14-1 p
14. --- Ik bewonderde de \VUgras{transatlantische schepen} \textsuperscript{(34)} die de rede
uitvaren, met hun \VUgras{vlaggen} \textsuperscript{(35)} in volle wind, geleid door een
behendige \VUgras{loods} \textsuperscript{(36)} die wonderwel het \VUgras{vaargeultje} weet te
volgen dat door \VUgras{boeien} \textsuperscript{(40)} en \VUgras{bakens} gemarkeerd is,
terwijl hij de \VUgras{lichters} \textsuperscript{(41)} vermijdt die zich moeizaam met hun
enige vierkante zeil laten sturen. Ik onderscheidde op het
\VUgras{voorschip} \textsuperscript{(37)} --- de boeg --- de \VUgras{hutten}
\textsuperscript{(39)} van de tweede klasse, en op het \VUgras{achterschip}
\textsuperscript{(38)} --- de spiegel --- de salons voor de passagiers van de eerste klasse.

%%K t10-15-1 p
15. --- Soms woonde ik ook het laden bij van een \VUgras{vrachtschuit} \textsuperscript{(42)}
waarin zojuist de goederen van het \VUgras{pakhuis} \textsuperscript{(43)} en van het
\VUgras{zeestation} \textsuperscript{(44)} opgestapeld werden, aangevoerd door de talrijke
treinen van de
\VUgras{kadelijn} \textsuperscript{(45)}.

%%K t10-tit-6 sub
\VUsaut{3.0mm}
\VUpk{10.2pt}{40}{Het roeien voor plezier.}
\VUsaut{3.0mm}

%%K t10-16-1 p
16. --- Dikwijls roeide ik in het \VUgras{vlotdok} \textsuperscript{(53)}, waarin de
\VUgras{boten} \textsuperscript{(46)} komen schuilen, en het was een vreugde voor mij de
\VUgras{riem} \textsuperscript{(47)} te hanteren. Ik klom op het \VUgras{dek}
\textsuperscript{(48)} van een \VUgras{zeilboot} \textsuperscript{(49)}, ik klom, langs
\VUgras{touwen} \textsuperscript{(52)}, in de \VUgras{mast} \textsuperscript{(51)} tot in de
\VUgras{ra's} \textsuperscript{(50)}, daarna herhaalde ik mijn tocht \VUgras{in een jol}
\textsuperscript{(54)} tussen zware \VUgras{vissersboten} \textsuperscript{(55)}, waar de
\VUgras{netten} \textsuperscript{(56)} drogen.

%%K t10-17-1 p
17. --- Op de \VUgras{kade} \textsuperscript{(58)} defileerden voor mij de meest uiteenlopende
\VUgras{goederen} \textsuperscript{(59)}, in \VUgras{kisten} \textsuperscript{(60)} of in
\VUgras{balen} \textsuperscript{(61)}. Zij vormden de \VUgras{lading} \textsuperscript{(63)}
van de \VUgras{lichters}, en krachtige
\VUgras{kranen} \textsuperscript{(62)} tilden ze op en zetten ze op
\VUgras{vrachtwagens} \textsuperscript{(64)} terwijl de \VUgras{voerlieden} ze
aangaven en de rechten in het \VUgras{douanekantoor} \textsuperscript{(65)} betaalden.

%%K t10-18-1 p
18. --- Ten slotte, toen ik bij de \VUgras{draaibrug} \textsuperscript{(66)} of bij de
\VUgras{sluis} \textsuperscript{(69)} kwam, langs de
\VUgras{baggermachine} \textsuperscript{(68)} die het \VUgras{kanaal}
\textsuperscript{(67)} schoonmaakt, ging ik op de kade aan wal bij de
\VUgras{seinpaal} \textsuperscript{(70)}.

%%K t10-19-1 p
19. --- Aan de andere kant van die kade komen wel de \VUgras{sloepen} \textsuperscript{(71)}
aanleggen die de officieren van het eskader aan land brengen. Het is een bewonderenswaardig
schouwspel de matrozen in de maat hun riemen te zien opheffen, terwijl de boot die door de
\VUgras{golf} \textsuperscript{(72)} gedragen wordt, gemakkelijk bij de trap van de
aanlegsteiger komt. Op die plaats kan men beter het \VUgras{getij} en de beweging van
\VUgras{eb} en \VUgras{vloed} op het \VUgras{strand} \textsuperscript{(73)} waarnemen.

%%K t10-tit-7 sub
\VUsaut{3.0mm}
\VUpk{10.2pt}{40}{De sociëteit.}
\VUsaut{3.0mm}

%%K t10-20-1 p
20. --- Om naar huis terug te keren gebruikte ik de \VUgras{elektrische tram}
\textsuperscript{(74)} waarvan de \VUgras{halte} \textsuperscript{(75)} bij de
\VUgras{paal} is die voor de sociëteit van de officieren staat.

%%K t10-20-2 p
Vanaf het \VUgras{balkon} \textsuperscript{(76)} van de sociëteit, versierd met
\VUgras{ankers} \textsuperscript{(77)}, kanonnen en kogels, zag ik dikwijls een
schitterende verzameling zeeofficieren: \VUgras{luitenants} \textsuperscript{(78)}, met hun
degen, \VUgras{kapiteins van de artillerie} \textsuperscript{(79)} en van de
\VUgras{infanterie} \textsuperscript{(80)} met hun \VUgras{sabel}. Achter hen stond een
\VUgras{matroos} \textsuperscript{(81)} die hun een \VUgras{verrekijker} bracht, wanneer zij
wilden onderscheiden wat er in de verte gebeurt.

%%K t10-21-1 p
21. --- Ik zag ook jonge \VUgras{luitenants-ter-zee} \textsuperscript{(82)} die de bevelen van
hun \VUgras{commandant} \textsuperscript{(83)} of van de
\VUgras{admiraal} \textsuperscript{(84)} ontvingen, terwijl een
\VUgras{kwartiermeester} \textsuperscript{(85)} wachtte om ze naar het schip te
brengen.

%%K t10-22-1 p
22. --- Van dat balkon is het uitzicht prachtig: men overziet het hele strand met zijn
\VUgras{tenten} \textsuperscript{(86)} en \VUgras{badkoetsjes} \textsuperscript{(87)}, en men
ziet, op het uur van het bad, de
\VUgras{baders} \textsuperscript{(88)} die vrolijk dartelen voor het
\VUgras{casino} \textsuperscript{(89)}.

%%K t10-23-1 p
23. --- Aan het uiteinde van het strand vormt de kust een klein rotsachtig
\VUgras{schiereiland} \textsuperscript{(91)}, waar de \VUgras{branding}
\textsuperscript{(92)} zich komt breken en waarop een \VUgras{vuurtoren}
\textsuperscript{(93)} gebouwd is, die 's nachts de weg aan de zeevarenden wijst. Dat
schiereiland is met het land alleen verbonden door een erg smalle
\VUgras{landengte} \textsuperscript{(90)}.

%%K t10-tit-8 sub
\VUsaut{3.0mm}
\VUpk{10.2pt}{40}{Algemeen gezicht op de kusten.}
\VUsaut{3.0mm}

%%K t10-24-1 p
24. --- De \VUgras{steile kust} \textsuperscript{(94)} strekt zich uit, opengereten door de
zee, die er grillig een reeks \VUgras{baaien} \textsuperscript{(95)} en \VUgras{voorgebergten}
(kapen) in tekent, waardoor zij heel schilderachtig wordt.

%%K t10-24-2 p
Op het \VUgras{plateau} \textsuperscript{(96)}, dat de \VUgras{(zee)engte}
\textsuperscript{(98)} beheerst, staan een heel belangrijk \VUgras{fort}
\textsuperscript{(97)} en enkele « villa's ».

%%K t10-25-1 p
25. --- Die engte scheidt van het vasteland een heel gevaarlijk
\VUgras{eiland} \textsuperscript{(99)}, omringd door \VUgras{klippen}
\textsuperscript{(100)} en \VUgras{branding}, waar boten en zelfs grote schepen soms stranden.
Ondanks de snelheid van de hulp, en de vlugge aankomst van de
\VUgras{reddingsboten}, zijn die \VUgras{schipbreuken} verschrikkelijk, en,
tijdens de \VUgras{stormen} van de nachtevening, zijn verscheidene schepen met man en goederen
vergaan.

%%K t10-25-2 p
Ondanks die nabijheid wordt \VUgras{de haven} veel door de zeelieden bezocht.
\VUsaut{6.0mm}
\VUfilet{22mm}
